\documentclass[11pt]{article}
\usepackage[a4paper,margin=24mm,headheight=15pt]{geometry}
\usepackage[T1]{fontenc}
\usepackage{lmodern,microtype,amsmath,amssymb,amsthm,booktabs,longtable,array,enumitem,fancyhdr,xurl}
\usepackage[hidelinks]{hyperref}
\newtheorem{theorem}{Theorem}[section]
\newtheorem{lemma}[theorem]{Lemma}
\newcommand{\Con}{\operatorname{Con}}
\newcommand{\file}[1]{{\small\path{#1}}}
\setlength{\parindent}{0pt}\setlength{\parskip}{5pt}
\setlist{nosep,leftmargin=1.6em,topsep=3pt}
\pagestyle{fancy}\fancyhf{}\fancyhead[L]{\small Two zero-spare command compositions at 215 states}\fancyhead[R]{\small 21 October 2026}\fancyfoot[R]{\thepage}
\hypersetup{pdftitle={Two zero-spare command compositions at 215 states},pdfauthor={Research contribution for the ProveIt project}}
\begin{document}
\begin{center}
{\LARGE\bfseries Two zero-spare command compositions\par at 215 states}\par\vspace{8pt}
{\large Finite removal of unsafe bootstrap data and protection of an unsafe tail}\par\vspace{5pt}
{Research contribution for Vladimir Reshetnikov's ProveIt project}\par
{21 October 2026 UTC}
\end{center}
\begin{abstract}
We construct two materially different binary, single-tape, blank-start machines with 215 working states. Neither is a 214-state result. Both combine the recently published zero-valued spare allocator with a different command language and a new initialization proof. The numerical calculus and its ZFC calibration are explicitly inherited.

The first machine uses command order NIP/MP/DROP. Its first complete source frame is \emph{unsafe}: finite initialization temporarily places the falsity code in both active theorem registers. Rather than asserting an invalid invariant there, we verify a further finite prefix. Failed numerical detachment removes the unsafe data, and after 570,263 binary transitions the machine reaches an ordinary safe root with next-candidate index 5 and additive counter 1. An arbitrary-counter calibration theorem then applies.

The second uses the three-command transfer-pop language. After 73,172 transitions it reaches a rigorously matched interior AX frame. Completing that AX leaves eight theorem heads above a generalization code not established to be a theorem. We do not call that code a theorem: an unbounded credit theorem bounds all future stack debt by three, permanently protecting the arbitrary suffix below at least five theorem heads. The two distinct initialization arguments are accompanied by full arbitrary-natural primitive runtime proofs and exact literal tables. These are ordinary mathematical proofs formalizable in ZFC, with executable certificates, not new Lean/Rocq kernel endpoints or claims of minimum state count.
\end{abstract}
\setcounter{tocdepth}{1}\tableofcontents
\section{Exact results and construction boundaries}
Each machine is deterministic, has one bi-infinite binary tape, and moves left or right at each transition. There are 215 working states, 430 transition rows, and a separate halt label. No nonblank tape is supplied externally. The claimed conclusion for each literal table $M$ is
\[
\mathrm{ZFC}\vdash\bigl(\operatorname{Halts}(M,\mathrm{blank})
\leftrightarrow\exists p\,\operatorname{Proof}_{\mathrm{ZFC}}(p,\bot)\bigr).
\]
Equivalently, blank nonhalting is provably equivalent in ZFC to $\Con(\mathrm{ZFC})$. This does not prove actual nonhalting, consistency of ZFC, or a lower bound on the required states.

\begin{center}\small
\begin{tabular}{lrr}\toprule
& NIP/MP/DROP & Transfer-pop\\\midrule
Working states &215&215\\
Binary transition rows &430&430\\
Primitive source occurrences &4,751&4,775\\
Distinct source subtrees &202&201\\
Guarded primitive cases &6,176&6,209\\
Symbolic rule applications &1,697,188&1,611,489\\
Minimum affine duration constant &6&7\\
Literal prefix used in proof &570,263&73,172\\
Prefix endpoint &Safe source root&Interior AX frame\\\bottomrule
\end{tabular}
\end{center}
Source DAG nodes are not physical states: the second source has 201 nodes but its machine has 215 states. Neither an initialized 214-state runtime nor a rejected 214-state table is substituted for these blank-start machines.

\subsection{A small table change with a nontrivial proof obligation}
For each construction, begin with its own pinned public 215-state parent, retain its source, change native row $198/0$ from $1L40$ to $0L40$, and designate native state 74 as initial rather than 0. All other rows are unchanged. The parent sources differ: the first is the previously uncalibrated order-201 command allocation; the second is the calibrated transfer-pop source. Their parents and exact recipes are included separately.

The changed row makes every source increment append a spare word
$B(0)=01$ rather than $B(1)=011$, where $B(v)=0\,1^{v+1}$. From matched frames, the appended spare region after $k$ increments has $2k$ rather than $3k$ cells. This saves $k$ cells in that region, not necessarily $k$ cells of total space at matched binary times. No general running-time improvement is asserted.

A different start state alone would not justify material novelty. Here the physical backend changes, and each combination requires a new source-level initialization argument. The first repairs an order that the public packet retained only as research. The second combines a different three-command language with zero spares and an unsafe-cache bootstrap that the previous unit-spare proof did not cover.

The number of one-writing rows falls by one relative to each unit-spare parent, excluding state-renaming equivalence even with global tape reflection. Further exact rooted graph comparisons against pinned public alternatives are included. These are bounded comparison claims, not assertions of global priority.

\subsection{Pinned tables}
\noindent\begin{minipage}{\linewidth}Primary native SHA-256:
\begin{quote}\small\path{ef49b414d3d8a08050654d225d4b6ec2972fc73d1bed2f277b5cfded88bdc101}\end{quote}\end{minipage}\par
\noindent\begin{minipage}{\linewidth}Primary canonical SHA-256:
\begin{quote}\small\path{21218bcf54b3feec8813b8266955f05de4f7f07c9a176ab830199f6747628f67}\end{quote}\end{minipage}\par
\noindent\begin{minipage}{\linewidth}Secondary native SHA-256:
\begin{quote}\small\path{29a6ed4a594d7f974cf267a9e6b665af34c9afd8dd7abbbecf0b538ecab151e0}\end{quote}\end{minipage}\par
\noindent\begin{minipage}{\linewidth}Secondary canonical SHA-256:
\begin{quote}\small\path{60eae1c30d214f71c65a207087dc716ba4cdad2a1a62c51e9e107434a0ceb4f1}\end{quote}\end{minipage}\par
Native tables start at 74 and halt at 254; their canonical tables start at 0 and halt at 215. Exact bijections are supplied. File hashes bind the literal serialization as well as the transition content.

\section{Representation and the zero-spare simulation}
The logical registers are ordered $(N,T,W,A,D,H,X,Y,L)$: next candidate, current theorem, encoded stack, selector residue, additive counter, theorem cache, two scratches, and proof-list remainder. The physical tape stores a fixed-width binary source address, a padding gap, unary primary registers, and a nonempty finite list of spare registers. The spare values need not be theorem codes and have no logical role.

The inherited symbolic checker is parameterized by the appended spare value. Here that parameter is zero, and the changed append block is rederived against each exact table. The domain still includes every natural primary register tuple and padding gap, arbitrary nonempty finite lists of natural spare values, and the complete zero exterior. In particular, mixed zero/one spare lists reached during bootstrap are permitted.

Unary scan and spare-list induction establish unbounded primitive contracts. Every guarded occurrence has a positive affine binary duration, so repeated source execution corresponds to genuine progress. A finite search cutoff for finding a symbolic derivation is not a bound on the natural registers quantified by a successful derivation. Conversely, a timeout would not prove a contract false. Full successful inventories are supplied separately for both tables.

The logical arguments below are independent of spare values. They nevertheless cannot replace the binary proof: the altered append row, source routing, each guard, and each exact initial configuration all require their own verified bindings.

\section{The numerical calculus and its consistency calibration}
\label{sec:native-calibration}

The present construction uses the dominated-parameter calculus supplied in
\file{K218_Dominated_Rho8_2026-10-20.zip}.  This is a recent input to the
repository, distinct from its primary, unchanged-calculus K218 table.
Its logical theorem is inherited through the 217-state parent. Both present constructions retain that numerical calculus. Their new contributions are the zero-spare command compositions and exact bootstrap calibration below.  This distinction is essential for identifying
which argument supplies each part of the result.

\subsection{Raw numerical operations}

For natural numbers put
\[
 P(a,b)=\frac{(a+b)(a+b+1)}2+b,
 \qquad I(a,b)=4P(a,b),\qquad a\to b:=I(a,b).
\]
The remaining constructors are
\[
 R(a,b)=4P(a,b)+1,\quad Q_a b=4P(a,b)+2,
 \quad E(a,b)=4P(a,b)+3.
\]
Raw truth is $T=0$ and the refutation target is $F=1=R(0,0)$.
Write $NX=I(X,F)$ and $\exists_a X=N(Q_a(NX))$.  In particular,
$I(0,0)=0$: zero is simultaneously the designated initial theorem and
an implication code.  All quantified schema parameters below range over
\emph{every natural number}.  No test that a number decodes as a conventional
formula is introduced.

Let $\mathcal C$ be the accepted native calculus, called \texttt{C222} in
the inherited interfaces and used by K219.  It has truth, numerical modus
ponens, and theorem-only generalization:
\[
 \frac{\mathcal C\vdash X\quad \mathcal C\vdash I(X,Y)}
      {\mathcal C\vdash Y},
 \qquad
 \frac{\mathcal C\vdash X}{\mathcal C\vdash Q_aX}.
\]
Its eight logical rows comprise the following seven retained families
and the unrestricted family $L$ displayed below:
\begin{align*}
 \rho_1(a,b)&=\exists_a E(a,b),\\
 \rho_2(a,b,c)&=Q_a(b\to c)\to(Q_ab\to Q_ac),\\
 \rho_3(a,b,c)&=E(a,b)\to(R(c,a)\to R(c,b)),\\
 \rho_4(a,b,c)&=Q_aQ_bc\to Q_bQ_ac,\\
 \rho_5(a,b,c)&=E(c,a)\to(R(c,b)\to R(a,b)),\\
 \rho_6(a,b)&=(\exists_a Q_ab)\to b,\\
 G(a,b,c)&=E(c,a)\to\bigl(E(c,b)\to Q_{R(a,b)}E(a,b)\bigr),\\
 L(A,B,C,U)&=((C\to U)\to A)\to((A\to C)\to(B\to C)).
\end{align*}

The accepted set-theoretic information remains in the literal gated seed
bundle.  To specify this inherited object, define
\begin{align*}
 B(b,k,m)&=Q_b\bigl(m\to NQ_b(m\to NR(b,k))\bigr),\\
 S(a,b,k)&=Q_k\bigl(R(k,a)\to R(k,b)\bigr),\\
 C(a,b,k,m)&=Q_b\bigl(R(b,k)\to
           NQ_a(m\to NQ_a(S(a,b,k)\to B(b,k,m)))\bigr),\\
 J(a,b,c)&=Q_c\bigl(C(a,b,c,c)\to NB(b,c,c)\bigr),\\
 V(a,b,c)&=R(a,b)\to Q_cR(a,b),\\
 \operatorname{Bundle}(a,b,c)&=N\bigl(V(a,b,c)\to J(a,b,c)\bigr).
\end{align*}
A bundle is an axiom when $c\ne a$ and $c\ne b$.  Both occurrences of
$c$ as matrix code and private label are retained literally.  We neither
remove these gates nor replace the bundle by an informal set-theoretic
analogue.  The source returns truth instead when a gate fails.

The established native calibration supplies finite-refutation
transformations between $\mathcal C\vdash F$ and ZFC inconsistency.
That inherited result is a premise of the present construction.  It is
supported by the repository's native-calculus developments and their
connection to the compact set theory underlying K230; it is not proved
anew by a transition-table simulation.

For provenance, the dominated package pins its K219 inputs to ProveIt
revision \path{f8ec6158f3bdc885139ee395ae34fe44c3251b9f} and includes
its exact native interfaces and finite logic certificate.  Earlier proof-assistant endpoints concern their own exact tables. The
native-calculus bridges, dominated theorem and present source/runtime
arguments must all be included to reach the new table. No new exact-215
proof-assistant endpoint or internal PA proof object is asserted here.

\subsection{Why a dominated fourth parameter is sufficient}

Define $\mathcal D$ by retaining all of the preceding clauses, except
that its $L(A,B,C,U)$ axiom requires
\[
                         U\le A.                 \tag{D}
\]
This is an ordinary inequality between natural codes, not a syntactic
freshness test and not an inequality of semantic truth values.

\begin{theorem}[Inherited dominated-parameter theorem]
For every natural code $n$,
\[
                  \mathcal D\vdash n
             \quad\Longleftrightarrow\quad
                  \mathcal C\vdash n.
\]
The proof is uniform in arbitrary native codes, including zero and
all degenerate constructor coincidences.
\end{theorem}

The easy direction follows because every dominated instance is already
an unrestricted instance.  For the converse, the inherited finite
certificate derives the entire unrestricted $L$ schema from truth,
modus ponens, and dominated instances alone.  Here are the arithmetic
and proof-theoretic facts that make that certificate meaningful.

First, $P(a,b)\ge a,b$ for all naturals.  If $a+b=s\ge1$, then
$s(s+1)/2\ge s$, and the case $s=0$ is immediate.  Consequently
$I(a,b)\ge a,b$.  If a schematic term $U$ is literally a subterm of an
implication term $A$, then $U\le A$ under every numerical assignment.
The same conclusion holds whenever $U=T$.  This remains true when some
implications evaluate to zero.  The finite checker uses precisely these
sufficient conditions for all dominated leaves.

Second, every certificate node denotes a universally quantified
\emph{numerical} statement.  Replacing one of its metavariables by an
implication term is therefore universal instantiation.  It does not
assume a general substitution rule inside the object calculus, and it
does not rename the numerical labels of a previously decoded atom.
At an MP node, equality of the instantiated antecedent and minor is
checked exactly.  The sole native normalization is $I(T,T)=T$.

The unrestricted-row certificate has the following exact inventory:
\begin{center}
\begin{tabular}{lr}
\hline
Item & Count\\
\hline
Proof nodes & 3,728\\
Shared term nodes & 5,160\\
Dominated-$L$ leaves & 1,260\\
Truth leaves & 1\\
MP steps with explicit substitutions & 55\\
Direct universal-instantiation steps & 410\\
Direct MP steps & 2,002\\
Maximum checked subterm depth & 4\\
\hline
\end{tabular}
\end{center}
These are finite proof objects, not samples of numerical assignments.
The implication-subterm inequality proves each axiom side condition for
all assignments; induction through the checked proof proves the
schematic result for all assignments.

The certificate first obtains identity, weakening $K$, Peirce's law,
composition, and distribution $S$:
\begin{align*}
 K&:p\to(q\to p),\\
 \mathrm P&:((p\to q)\to p)\to p,\\
 \mathrm B&:(p\to q)\to((q\to r)\to(p\to r)),\\
 S&:(p\to(q\to r))\to((p\to q)\to(p\to r)).
\end{align*}
A key intermediate theorem is
\[
 (p\to q)\to\bigl(((p\to r)\to p)\to q\bigr).
\]
The successful inherited proof of that theorem adds the dominated
context $A=((z\to U)\to v)\to w$; the fourth argument lies three
implication edges inside the first.  The rest of the fixed derivation
reaches unrestricted $L$.  The complete certificate and its checkers
are included with the inherited dependency, so the large intermediate
formulas need not be expanded into the body of this article.

For clarity, the final passage from $K,S,\mathrm P$ to $L$ has a short
ordinary explanation.  Temporarily assume
$H=(C\to U)\to A$ and $J=A\to C$.  Under the additional assumption
$C\to U$, two MPs give $A$ and then $C$.  The deduction transformation
justified by $K,S$ gives $(C\to U)\to C$.  Peirce's law at $C,U$
gives $C$.  Weakening gives $B\to C$, and discharging $J,H$ yields
exactly $L(A,B,C,U)$.  No generalization is used under these temporary
assumptions.  The supplied certificate expands this deduction argument
into explicit schema instances and MP steps.

Replacing each unrestricted-$L$ leaf of a finite $\mathcal C$ proof by
this fixed derived proof gives a finite $\mathcal D$ proof of the same
code.  All other axioms and both rules are retained literally, proving
the theorem.  Combining it with the inherited native calibration gives
\[
       \mathcal D\vdash1
       \quad\Longleftrightarrow\quad
       \text{there is a finite ZFC refutation}.       \tag{C1}
\]


\section{Primary construction: removing unsafe data during a finite prefix}
\subsection{The command language and the unsafe first root}
Let $P$ be Cantor pairing and $I(a,b)=4P(a,b)$. Write the retained tuple as $(t,w,d,h)$. Order 201 means
\[
\mathsf{NIP}:0,\qquad \mathsf{MP}:1,\qquad
\mathsf{DROP}:2,\qquad \mathsf{AX}(c;x,y):c+3.
\]
The non-AX field lists are $[0]$, $[1,x]$, and $[2]$ respectively; AX fields are $[c+3,y,x]$. The primitive bodies and their all-natural command meanings are inherited. Independent control resolvers check the selector partitions, field arities, and normal returns for this exact source.

NIP tests $h=1$ \emph{before} popping. If it does not halt and $(x,v)=P^{-1}(w)$, it returns $(t,v,d,x)$. DROP returns $(h,P(t,w),d,h)$. MP$(x)$ puts $y=d+x$ and returns
\[
\begin{cases}(y,P(0,w),y,0),&I(t,y)=h,\\(0,w,y,0),&I(t,y)\ne h.\end{cases}
\]
AX puts $U=d,A=d+x,B=h+y,C=c$, returns the gated bundle as top, resets $d=h=0$, and pushes the seven logical rows, the dominated row, $Q_Uh$, and the old top. The bound $U\le A$ is automatic for every natural input.

Exact blank execution first reaches a complete zero-address source frame after 118,411 transitions with
\[
(N,T,W,A,D,H,X,Y,L)=(2,0,2,0,0,0,0,0,0).
\]
This is \emph{not} a theorem-stack seed: $2=P(0,1)$ contains truth above an unsafe tail. No consistency invariant is applied at that point.

\subsection{Four commands that discharge the unsafe suffix}
The exact candidate packing determines the next commands, without assumptions on ZFC. Since $2=P(0,1)$ and subsequent positive remainders are decremented before decoding, candidate 2 executes NIP,NIP. Candidate $3=P(2,0)$ executes DROP. Candidate $4=P(1,1)$ executes MP$(1)$ because $1=P(1,0)$ supplies the payload.

\begin{center}\small
\begin{tabular}{lrrrr}\toprule
Boundary & $T$&$W$&$D$&$H$\\\midrule
Initial complete frame&0&2&0&0\\
After first NIP&0&1&0&0\\
After second NIP&0&0&0&1\\
After DROP&1&0&0&1\\
After MP$(1)$&0&0&1&0\\\bottomrule
\end{tabular}
\end{center}
Neither NIP halts: both test cache zero, and the second only \emph{then} loads falsity. DROP has no halt test. The final MP fails because $I(1,1)=16\ne1$, clearing both active formula registers. It leaves only the harmless additive counter value 1.

An independently implemented interval-control interpreter and the arithmetic control resolver agree on all source edges. Both reproduce the 435-primitive finite segment, with root visits after 0, 63, 89, 166, and 435 primitives. No halting primitive occurs. The terminal root is
\[
                         (5,0,0,0,1,0,0,0,0).
\]
A counterfactual theoremhood assertion during the earlier segment would be wrong. A verified finite nonhalting prefix needs no such invariant; the invariant begins only at its safe endpoint.

\subsection{A complete literal binary prefix}
The proof does not rely solely on the source interpreter for initialization. Independent native and canonical simulation replays all 570,263 binary transitions from blank tape, never halting, and reaches native state 84, head $-9$, gap 131, with the displayed safe tuple and 176 natural spare words.

The full sparse tape is reconstructed and compared. Its minimum one-cell is $-10=h-1$, and its next one is $186=h+64+131$. Hence all 64 source-address cells are zero. This is a genuine root, not an interior dispatcher frame. Every primary word, every mixed spare value, and the entire zero exterior are checked.

\subsection{Strong proof compilation and arbitrary initial counters}
The inherited strong compiler supplies, for each numerical theorem $F$, a finite named-command block $\mathcal B_F$ which from every $(t,w,0,0)$ either halts or returns $(F,w,0,0)$. It quantifies over arbitrary incoming top and stack. Command renumbering transports it to order 201 by changing wire tags, not arithmetic meanings.

For clarity, its principal blocks can be given explicitly. Put
\[
\mathsf{Discard}=\mathsf{MP}(1);\mathsf{AX}(1;0,0);\mathsf{NIP}^{10},
\quad J(u)=\begin{cases}[],&u=0,\\{}[\mathsf{MP}(u)],&u>0.\end{cases}
\]
From $D=H=0$, Discard returns truth with the old stack and zero counter/cache. From a zero top/counter/cache, $J(u)$ sets the counter to $u$; omitting MP when $u=0$ avoids its zero-alias successful push. Truth uses Discard. A retained row $R_j(a,b,c)$ uses
\[
\mathsf{Discard};\mathsf{AX}(c;a,b);\mathsf{NIP}^j;
\mathsf{DROP};\mathsf{NIP}^{11-j}.
\]
The dominated row with $u\le a$ uses
\[
\mathsf{Discard};J(u);\mathsf{AX}(c;a-u,b);
\mathsf{NIP}^8;\mathsf{DROP};\mathsf{NIP}^{3}.
\]
The bundle uses $\mathsf{Discard};\mathsf{AX}(c;a,b);\mathsf{NIP}^{10}$ when its gates hold. All generated row entries are nonfalsity implication codes and the generalized cache has tag 2, so these selection suffixes restore the underlying stack exactly.

For generalization use
\[
\mathcal B_X;\mathsf{DROP};J(i);\mathsf{NIP};\mathsf{AX}(i;0,0);
\mathsf{NIP}^{9};\mathsf{DROP};\mathsf{NIP};\mathsf{NIP}.
\]
For modus ponens from $X$ and $I(X,Y)$ use
\begin{multline*}
\mathcal B_X;\mathsf{DROP};\mathcal B_{I(X,Y)};
\mathsf{NIP};\mathsf{DROP};\mathsf{NIP};\mathsf{MP}(Y);\mathsf{AX}(Y;0,0);\\
\mathsf{NIP}^{10};\mathsf{DROP};\mathsf{NIP};\mathsf{NIP}.
\end{multline*}
The AX reset has equal selector/counter parameter and therefore zero gate; the last selection restores $Y$ and the old stack, or halts earlier on falsity. These exact contracts, including the zero implication alias, prove the compiler by induction on a finite numerical derivation.

\begin{theorem}[Safe root with arbitrary counter]
At a genuine root with $N=n,D=d$ arbitrary naturals, $A=X=Y=L=0$, and theorem-valued top, cache, and stack heads, source halting is equivalent to numerical inconsistency.
\end{theorem}
\begin{proof}
Soundness permits every $d$: AX uses $U=d\le d+x=A$, and all other commands retain their theoremhood arguments. Only cached falsity can trigger halt.

For completeness, append DROP,NIP,NIP to $\mathcal B_1$. Prefix by $\mathsf{AX}(0;1,y)$ with $y\ge1$. It resets counter/cache from every incoming tuple; its resulting top and stack may be arbitrary because the strong block permits them. If $R$ is the successor-based tail encoding of the fixed refutation suffix, the exact first-candidate number is
\[
                         n_y=P(3,P(y,P(1,R)))\ge y.
\]
There is no outer successor on a first candidate. Choose $y\ge\max(1,n)$. Every earlier nonhalting candidate drains, because arithmetic terminates and subsequent positive remainders strictly decrease before tail projection. The scheduler thus reaches a successful candidate unless it halts earlier. No bound on an unknown incoming counter is assumed.
\end{proof}
The verified endpoint has $n=5,d=1$ and otherwise zero theorem data, so this theorem applies directly. No protected-tail debt theorem is needed for this primary construction.

\section{Secondary construction: protecting an unsafe generalization tail}
\subsection{The transfer-pop language}
The second source has three commands: composite transfer-pop $S$ at tag 0, MP at tag 1, and AX at tags $c+2$, with field lists $[0]$, $[1,x]$, and $[c+2,y,x]$. Its MP and AX meanings agree with the preceding source. The actual $S$ is old DROP; old NIP; old NIP. Therefore it halts if \emph{either} the incoming top or cache is 1; otherwise, writing $P^{-1}(w)=(x,v)$, its normal result is
\[
                             S(t,w,d,h)=(h,v,d,x).
\]
A cache-only halt description would be incomplete. The imported compiler first analyzes that idealized weaker guard and then lifts it to the actual composite by early-halting refinement. On safe actual executions, both extra tests concern theorem codes, preserving soundness; in the completeness contract an earlier halt is already success.

The initial-state/append-row change leaves all 4,775 source primitives unchanged. The complete AX subtree comprises 4,565 leaves; an independent comparison verifies both their literal operations and their relative paths against the parent.

\subsection{An exact interior source boundary, not an invented root}
After 73,172 binary transitions from blank tape, the machine is at native dispatcher 84, head $-2$, gap zero, with 65 natural spare words. The logical tuple is
\[
                        (0,6,0,0,1,1,10,1,0).
\]
Its source ordinal is 119, with address \texttt{11010011000110110110}. This address is \emph{not} zero. Both native and canonical simulators compare every transition and the whole terminal tape; no claim of a root seed is made at this point.

A separate structured-source interpreter begins at AX entry ordinal 25 with
\[
                         (N,T,W,A,D,H,X,Y,L)=(0,0,0,0,1,1,0,0,0)
\]
and reaches exactly the observed ordinal 119 and entire register store after 177 primitives. The remaining AX execution is total on these natural inputs. Positive unbounded primitive simulation therefore transports the binary machine through the same remaining source suffix. Literal execution of its potentially enormous output codes is not required.

\subsection{Eight safe heads, one deliberately untrusted head}
The virtual complete AX has $U=A=B=1,C=0$. Its gate is \emph{enabled}, because $0\ne1$ on both tests. Its output top is the accepted bundle $K(1,1,0)$, and it resets counter/cache to zero. Its stack is
\[
[R_1(1,1,0),\ldots,R_7(1,1,0),L(1,1,0,1),Q_1(1),0;0].
\]
The first eight heads are theorems: the seven retained rows and the dominated row with $U=A=1$. The incoming cache 1 has not been established to be a theorem. We therefore do \emph{not} claim that its generalization $Q_1(1)$ is a theorem.

Instead isolate an arbitrary suffix
\[
       J=P(Q_1(1),P(0,0))=P(18,0)=171.
\]
The represented stack consists of eight explicit theorem heads above $J$. The top is a theorem bundle and the cache is truth. The next-candidate counter remains $N=0$ and the proof-list remainder is zero. The following unbounded scheduler argument protects that suffix forever, without any theoremhood assumption on its contents.

\subsection{The infinite schedule has conservative debt three}
Write $f(n)$ for the right projection of Cantor pairing and $T_s=s(s+1)/2$. Let $W(n)$ be the finite syntactic command word of a freshly fetched candidate, ignoring possible semantic halts. Put $R(0)=[]$ and $R(r)=W(r-1)$ for $r>0$. If $P^{-1}(n)=(a,y)$, the exact source field arities give
\[
W(n)=\begin{cases}
S;R(y),&a=0,\\
\mathsf{MP};R(f(y)),&a=1,\\
\mathsf{AX};R(f(f(y))),&a\ge2.
\end{cases}
\]
These words terminate by strict descent of the recursive index. Each $R(r)$ has at most $r$ commands. Assign conservative stack balances $-1,0,+10$ to $S$, MP, AX. A successful MP adds a further theorem zero; ignoring it only weakens the protection bound.

Let $b(n)$ be the balance of $W(n)$, $B(k)=\sum_{n<k}b(n)$, $g(0)=0$, and $g(r)=b(r-1)$ for $r>0$. If a sequence $q$ has nonnegative initial partial sums, so does $q(f(0)),q(f(1)),\ldots$: the projection sequence concatenates $[0],[0,1],[0,1,2],\ldots$. The assertion follows by summing complete blocks and one partial block, and can be applied twice.

Use the shifted sequence $q(r)=g(r)+1$. Its first $k\ge1$ terms sum to $1+B(k-1)+(k-1)$. The exceptional candidates are
\[
                         W(0)=S,\qquad W(1)=\mathsf{MP},\qquad W(2)=SS,
\]
so $B(0),B(1),B(2),B(3)=0,-1,-1,-3$, and $B(k)+k\ge0$ through $k=3$.

Induct on the Cantor diagonal $s\ge2$. Its first $s-1$ candidates start with AX, followed by one MP candidate and one S candidate. For the first $k$ AX candidates the balance is
\[
10k+\sum_{y<k}g(f(f(y)))\ge9k,
\]
by the twice-applied prefix lemma. It uses only previously proved inequalities at indices below $T_s$, so the induction is not circular. Inside AX candidate $y$, its continuation loses at most $y$, giving lower bound $9y+10-y=8y+10$. The following MP candidate loses at most $s-1$, leaving at least $8(s-1)$. The final S and its continuation lose at most $1+s$, leaving
\[
                              8(s-1)-1-s=7s-9>0.
\]
Every command prefix of the diagonal is therefore nonnegative relative to its start. The diagonal boundary remains at least $-3$, and $B(k)+k\ge0$ continues to hold, completing the induction. The global conservative minimum from cutoff zero is exactly $-3$, attained in the exceptional first three candidates.

\begin{lemma}[Protection of the unsafe suffix]
Starting with eight theorem heads above $J$, theorem-valued top/cache, $N=L=0$, the arbitrary suffix $J$ is never exposed during a nonhalting execution. At least five explicit theorem heads remain above it at every command boundary.
\end{lemma}
\begin{proof}
Track the explicit wrapper length using the conservative balance. AX adds ten theorem heads once the active top/cache are safe; MP preserves or adds a theorem head. Before an $S$ pop the conservative count is at least one more than its post-pop count, which is at least five. Thus it removes a safe head. Inside composite $S$, DROP first pushes a temporary theorem top, the first NIP removes that temporary head, and the second removes precisely the safe old head accounted for by the balance. Its two halt tests also inspect only theorems.

The wrapper is an explicitly maintained finite list above a fixed tail, not a claim that its length is recoverable from the encoded integer. This handles $P(0,0)=0$: the exact identity $P^{-1}(P(x,v))=(x,v)$ still supports each push/pop. The induction applies to every finite actual prefix; defining a syntactic schedule while ignoring halts does not assume nonhalting.
\end{proof}
Consequently every actual halt certifies numerical falsity as a theorem. This supplies soundness despite the initial untrusted cache and buried code 18.

\subsection{Completeness is inherited through a strong arbitrary-stack compiler}
The transfer source has a complete named-command compiler, not merely a theoremhood invariant. Its strong contract is the same: from every $(t,w,0,0)$, a block for $F$ either halts or returns $(F,w,0,0)$. The parent packet contains the full derivation and its separate composite-guard review.

Two structural macros explain why removing primitive DROP does not lose completeness. Put $D_0=\mathsf{AX}(0;0,0)$. Then
\[
 \mathsf{PUSH}=D_0;S^9;\mathsf{MP}(0),\qquad
 \mathsf{POP}=D_0;S^{11}.
\]
For the ideal cache-guard transfer, PUSH returns $(0,P(t,w),0,0)$ on every input, and POP preserves a nonfalsity top while popping one old head and resetting the counter, or halts when the top is 1. For the actual composite, any additional earlier halt is allowed by the strong compiler contract. If no such halt occurs, the displayed normal outputs remain identical. This early-halting distinction is essential; PUSH is not asserted to be nonhalting on arbitrary unsafe inputs for the actual composite.

Axiom selection, theorem generalization, and numerical detachment are built from these macros with exact ten-entry stack cleanup. A compiled derivation of 1, followed by POP, halts on every retained top/stack after the counter and cache are reset. Prefix by $\mathsf{AX}(0;1,y)$ as before. Its first-candidate integer is now
\[
                              P(2,P(y,P(1,R)))\ge y,
\]
because AX begins at tag 2. These successful candidates are cofinal. Finite draining of earlier candidates proves completeness from cutoff zero. The arbitrary tail $J$ causes no obstruction: the compiler's input contract explicitly permits it, while the independent credit argument proves actual-run soundness.

The exact interior bootstrap, total AX suffix, credit theorem, strong compiler, and unbounded binary simulation therefore establish the blank-start consistency equivalence for the secondary table.

\section{Verification, reproduction, and limits}
Each claimed result has five separate components:
\begin{enumerate}
\item A total literal binary table with 215 working states, exact start and halt conventions, and a checked native/canonical bijection.
\item A precise source and its command-calculus calibration. The first uses a tag-permuted strong compiler; the second uses the transfer compiler and its composite-guard refinement.
\item Complete arbitrary-natural primitive contracts, with inductive scan/list lemmas and positive binary durations. The domains include unsafe natural formula codes, which is necessary for both bootstrap proofs.
\item An exact finite blank prefix checked against the whole tape, not just a few decoded registers. The first continues to a safe zero-address root; the second stops at a matched interior AX frame and uses a separately proved total suffix.
\item An all-natural soundness/completeness argument from the calibrated boundary, connected to the explicitly inherited numerical-calculus/ZFC theorem.
\end{enumerate}
Finite source, command, graph, and startup checks support their stated finite claims. They do not substitute for unbounded primitive or scheduler proofs. The mathematical proof engine is executable, but not a proof-assistant kernel. A fresh rerun of that engine is not a second independently designed proof algorithm. Independent parsers, native/dense tape replay, a separately implemented source-prefix check, and mathematical review of the credit proof are identified as such in the receipts.

The archive includes portable entry points for default reconstruction/binding checks and optional complete fresh runtime derivations, together with literal tables, source recipes, compressed lossless traces, proof texts, attribution, and SHA-256 manifests. Its README specifies the exact commands and distinguishes checking saved inventories from rederiving the symbolic proofs. Python assertions must be enabled; entry points reject optimized Python rather than silently skipping checks.

\subsection{What remains open at 214}
The investigated 214-state initialized core still fails blank initialization, including with zero-valued spares. Concrete all-start outward-ray certificates exclude those exact tables and starts, not all possible constructions. Larger source re-encoding and helper-allocation searches also found no certified 214-state blank machine. Timeouts and restricted allocation-model failures are not state lower bounds.

The present deliverable satisfies the alternative criterion of materially different 215-state constructions. It does not rename a failed 214-state table, count source nodes as states, supply a nonblank initial tape, or claim that a source-only compiler already gives a smaller binary machine.

\subsection{Provenance}
The repository is \url{https://github.com/VladimirReshetnikov/ProveIt}. The public zero-spare packet and command-family source were introduced in commit \file{cd596116875}; the transfer-pop packet is the one introduced through commits \file{7342ae8c96c} and \file{2c77edead6a}. The new work combines those inputs and proves the distinct bootstrap calibrations. The underlying numerical calculus, dominated-row certificate, arithmetic architecture, command compilers, credit theorem, and runtime induction checker retain their upstream attribution and licenses.

The primary numerical-calculus soundness boundary is reached only after the complete 570,263-step finite prefix. The secondary boundary contains only eight assumed theorem heads, not all ten AX outputs. These qualifications are part of the theorem, not merely implementation notes. Neither the repository's uncompiled endpoint drafts nor a literal-data audit is cited as a new kernel-checked endpoint for these tables.

\clearpage\appendix
\section{Complete canonical table: NIP/MP/DROP machine}
The working labels are $0,\ldots,214$, the initial state is 0, and $H$ is halt label 215 with no outgoing rows. Initially the tape is everywhere zero and the head is at coordinate 0. An entry $wDq$ writes $w$, moves one cell in direction $D$, and enters $q$.
\small\renewcommand{\arraystretch}{0.94}
\begin{longtable}{rll@{\qquad}rll}
\toprule State & Read 0 & Read 1 & State & Read 0 & Read 1\\\midrule
\endfirsthead
\toprule State & Read 0 & Read 1 & State & Read 0 & Read 1\\\midrule
\endhead
0 & $1R1$ & $1R65$ & 108 & $0R123$ & $1R9$\\
1 & $1R93$ & $1R1$ & 109 & $0R102$ & $1R124$\\
2 & $0L2$ & $1L3$ & 110 & $0R152$ & $1R190$\\
3 & $0L4$ & $1L3$ & 111 & $0L214$ & $0L96$\\
4 & $0L6$ & $1L3$ & 112 & $0R126$ & $1R105$\\
5 & $1R7$ & $0L5$ & 113 & $0R201$ & $1R29$\\
6 & $0L10$ & $1L3$ & 114 & $0R167$ & $1R184$\\
7 & $0R7$ & $0R11$ & 115 & $0R76$ & $1R137$\\
8 & $0R85$ & $1R194$ & 116 & $0R150$ & $1R148$\\
9 & $0R15$ & $1R14$ & 117 & $0R162$ & $1R205$\\
10 & $0L16$ & $1L3$ & 118 & $0L133$ & $0L2$\\
11 & $0L17$ & $0R11$ & 119 & $0R40$ & $1R74$\\
12 & $0L27$ & $1R138$ & 120 & $0R63$ & $1R134$\\
13 & $0R20$ & $1R21$ & 121 & $0R63$ & $1R135$\\
14 & $0L90$ & $1L32$ & 122 & $0R72$ & $0L17$\\
15 & $0L12$ & $0R23$ & 123 & $0R15$ & $1R134$\\
16 & $0L24$ & $1L3$ & 124 & $0R22$ & $1R136$\\
17 & $1L2$ & $0L2$ & 125 & $0R149$ & $1R181$\\
18 & $0R25$ & $0R33$ & 126 & $0R139$ & $1L21$\\
19 & $0R206$ & $1R205$ & 127 & $0R132$ & $1R183$\\
20 & $0R29$ & $1R30$ & 128 & $0R209$ & $1R174$\\
21 & $0R31$ & $1R32$ & 129 & $0R76$ & $1R179$\\
22 & $0R25$ & $0R61$ & 130 & $0R169$ & $1R95$\\
23 & $0R34$ & $1R23$ & 131 & $0R70$ & $1R200$\\
24 & $0L35$ & $1L3$ & 132 & $0R13$ & $1R143$\\
25 & $0L36$ & $0R1$ & 133 & $0R31$ & $1R71$\\
26 & $0R37$ & $1R26$ & 134 & $0R40$ & $1L126$\\
27 & $0R38$ & $1R39$ & 135 & $0R155$ & $0L17$\\
28 & $0R40$ & $0L41$ & 136 & $0R156$ & $0L17$\\
29 & $0R42$ & $1R43$ & 137 & $0R163$ & $1R130$\\
30 & $0R44$ & $1R45$ & 138 & $0R208$ & $0L214$\\
31 & $0R46$ & $1R47$ & 139 & $0R25$ & $0R49$\\
32 & $0R22$ & $1R48$ & 140 & $0R142$ & $1R183$\\
33 & $0R49$ & $1R33$ & 141 & $0R12$ & $1R158$\\
34 & $0R33$ & $1R34$ & 142 & $0R97$ & $1R161$\\
35 & $0L51$ & $1L3$ & 143 & $0R178$ & $1R168$\\
36 & $1R36$ & $1L2$ & 144 & $0R13$ & $1R161$\\
37 & $0R52$ & $1R37$ & 145 & $0R89$ & $1L168$\\
38 & $0R18$ & $1R202$ & 146 & $0R89$ & $1R160$\\
39 & $0R14$ & $1R54$ & 147 & $0L154$ & $1L165$\\
40 & $0L12$ & $0R37$ & 148 & $0R85$ & $1R89$\\
41 & $0L5$ & $1L55$ & 149 & $0R161$ & $1R13$\\
42 & $0R25$ & $0R26$ & 150 & $0R84$ & $1R145$\\
43 & $0R14$ & $0L17$ & 151 & $0R166$ & $1R170$\\
44 & $0R57$ & $1R58$ & 152 & $0L68$ & $1L143$\\
45 & $0R59$ & $1R60$ & 153 & $0R150$ & $1R207$\\
46 & $0R25$ & $0R56$ & 154 & $0R161$ & $1R116$\\
47 & $0R62$ & $1R63$ & 155 & $0R98$ & $1R98$\\
48 & $0R64$ & $0L17$ & 156 & $0R50$ & $0R33$\\
49 & $0R26$ & $1R49$ & 157 & $0R109$ & $1R27$\\
50 & $0L36$ & $0R65$ & 158 & $0R113$ & $1R159$\\
51 & $0L75$ & $1L3$ & 159 & $0L191$ & $1R97$\\
52 & $0R17$ & $1R52$ & 160 & $0R170$ & $1R182$\\
53 & $0R149$ & $1R145$ & 161 & $0R20$ & $1R173$\\
54 & $0R69$ & $1R70$ & 162 & $0R177$ & $1R167$\\
55 & $0L41$ & $1L55$ & 163 & $0R161$ & $1R97$\\
56 & $0R23$ & $1R56$ & 164 & $0R127$ & $1R166$\\
57 & $0L12$ & $0R52$ & 165 & $0R98$ & $1R164$\\
58 & $0R62$ & $0L17$ & 166 & $0L8$ & $1R144$\\
59 & $0R71$ & $1R72$ & 167 & $0L110$ & $1R177$\\
60 & $0R63$ & $1R73$ & 168 & $0R148$ & $1R176$\\
61 & $0R56$ & $1R61$ & 169 & $0R131$ & $1R212$\\
62 & $0R50$ & $0R26$ & 170 & $0R163$ & $1R145$\\
63 & $0L29$ & $1L74$ & 171 & $0R192$ & $1R167$\\
64 & $0L119$ & $1L31$ & 172 & $0R181$ & $1R181$\\
65 & $0R0$ & $1R65$ & 173 & $0R184$ & $1R124$\\
66 & $0R98$ & $1R152$ & 174 & $0R147$ & $1R167$\\
67 & $1R188$ & $1L55$ & 175 & $0R171$ & $1R88$\\
68 & $0R146$ & $1R147$ & 176 & $0R76$ & $1R62$\\
69 & $0R77$ & $1RH$ & 177 & $0R147$ & $1R89$\\
70 & $0L201$ & $1R32$ & 178 & $0R142$ & $1R85$\\
71 & $0R42$ & $1R64$ & 179 & $0R178$ & $1R76$\\
72 & $0L90$ & $1R79$ & 180 & $0R146$ & $1R175$\\
73 & $0R57$ & $1L80$ & 181 & $0R144$ & $1L151$\\
74 & $0R81$ & $1R82$ & 182 & $0R98$ & $1R53$\\
75 & $0R75$ & $1R83$ & 183 & $0R170$ & $1R53$\\
76 & $0L147$ & $1L143$ & 184 & $0R18$ & $1L133$\\
77 & $0R86$ & $1R87$ & 185 & $0R145$ & $1R180$\\
78 & $0R201$ & $1R142$ & 186 & $0R77$ & $1R118$\\
79 & $0R90$ & $1L134$ & 187 & $0R107$ & $1R126$\\
80 & $0L109$ & $1R58$ & 188 & $0L41$ & $1L111$\\
81 & $0R50$ & $0R34$ & 189 & $0R174$ & $1R161$\\
82 & $0R92$ & $0L43$ & 190 & $0R152$ & $1R117$\\
83 & $0R94$ & $1R157$ & 191 & $0R185$ & $1R186$\\
84 & $0R97$ & $1R13$ & 192 & $0R196$ & $1R100$\\
85 & $0R89$ & $1R98$ & 193 & $0R148$ & $1R203$\\
86 & $0R213$ & $1R22$ & 194 & $0R159$ & $1R195$\\
87 & $0R118$ & $1R101$ & 195 & $0R178$ & $1R189$\\
88 & $0R152$ & $1R114$ & 196 & $0R66$ & $1R8$\\
89 & $0R103$ & $1R104$ & 197 & $0R152$ & $1R129$\\
90 & $0R81$ & $1R80$ & 198 & $0R166$ & $1R191$\\
91 & $0R102$ & $1R99$ & 199 & $0R154$ & $1R125$\\
92 & $0R106$ & $1R107$ & 200 & $0R89$ & $1R78$\\
93 & $0R67$ & $0R1$ & 201 & $0R91$ & $1R89$\\
94 & $0R28$ & $1R108$ & 202 & $0R141$ & $0L214$\\
95 & $0R186$ & $0L214$ & 203 & $0R210$ & $1R110$\\
96 & $1L111$ & $1L96$ & 204 & $0R152$ & $1R172$\\
97 & $0R20$ & $1R112$ & 205 & $0R204$ & $1R162$\\
98 & $0R62$ & $1R62$ & 206 & $0R142$ & $1R140$\\
99 & $0R105$ & $1R97$ & 207 & $0R199$ & $1R19$\\
100 & $0R68$ & $1R193$ & 208 & $0R211$ & $1R13$\\
101 & $0R106$ & $0L214$ & 209 & $0R165$ & $1R197$\\
102 & $0R119$ & $1R120$ & 210 & $0R128$ & $1R153$\\
103 & $0R73$ & $1R60$ & 211 & $0R18$ & $1R115$\\
104 & $0R59$ & $1R121$ & 212 & $0R198$ & $1R187$\\
105 & $0R22$ & $1R122$ & 213 & $0R46$ & $1R101$\\
106 & $0R22$ & $0L11$ & 214 & $0L5$ & $0L214$\\
107 & $0R63$ & $1R43$ & & &\\
\bottomrule
\end{longtable}
\renewcommand{\arraystretch}{1}\normalsize

\clearpage
\section{Complete canonical table: transfer-pop machine}
The same conventions apply, with its own bijection from native labels. These tables are different machines, not two relabelings of one table.
\small\renewcommand{\arraystretch}{0.94}
\begin{longtable}{rll@{\qquad}rll}
\toprule State & Read 0 & Read 1 & State & Read 0 & Read 1\\\midrule
\endfirsthead
\toprule State & Read 0 & Read 1 & State & Read 0 & Read 1\\\midrule
\endhead
0 & $1R1$ & $1R64$ & 108 & $0R148$ & $1R183$\\
1 & $1R91$ & $1R1$ & 109 & $0L204$ & $0L94$\\
2 & $0L2$ & $1L3$ & 110 & $0R123$ & $1R103$\\
3 & $0L4$ & $1L3$ & 111 & $0R206$ & $1R180$\\
4 & $0L6$ & $1L3$ & 112 & $0R162$ & $1R178$\\
5 & $1R7$ & $0L5$ & 113 & $0R146$ & $1R144$\\
6 & $0L10$ & $1L3$ & 114 & $0R157$ & $1R197$\\
7 & $0R7$ & $0R11$ & 115 & $0L130$ & $0L2$\\
8 & $0R83$ & $1R187$ & 116 & $0R39$ & $1R73$\\
9 & $0R15$ & $1R14$ & 117 & $0R62$ & $1R131$\\
10 & $0L16$ & $1L3$ & 118 & $0R62$ & $1R132$\\
11 & $0L17$ & $0R11$ & 119 & $0R71$ & $0L17$\\
12 & $0L27$ & $1R128$ & 120 & $0R15$ & $1R131$\\
13 & $0R20$ & $1R21$ & 121 & $0R22$ & $1R133$\\
14 & $0L88$ & $1L32$ & 122 & $0R145$ & $1R175$\\
15 & $0L12$ & $0R23$ & 123 & $0R135$ & $1L21$\\
16 & $0L24$ & $1L3$ & 124 & $0R129$ & $1R177$\\
17 & $1L2$ & $0L2$ & 125 & $0R201$ & $1R168$\\
18 & $0R25$ & $0R33$ & 126 & $0R75$ & $1R173$\\
19 & $0R198$ & $1R197$ & 127 & $0R18$ & $1R174$\\
20 & $0R29$ & $1R30$ & 128 & $0R75$ & $1R134$\\
21 & $0R31$ & $1R32$ & 129 & $0R13$ & $1R139$\\
22 & $0R25$ & $0R60$ & 130 & $0R31$ & $1R70$\\
23 & $0R34$ & $1R23$ & 131 & $0R39$ & $1L123$\\
24 & $0L35$ & $1L3$ & 132 & $0R151$ & $0L17$\\
25 & $0L36$ & $0R1$ & 133 & $0R152$ & $0L17$\\
26 & $0R37$ & $1R26$ & 134 & $0R158$ & $1R137$\\
27 & $0R127$ & $1R38$ & 135 & $0R25$ & $0R48$\\
28 & $0R39$ & $0L40$ & 136 & $0R138$ & $1R177$\\
29 & $0R41$ & $1R42$ & 137 & $0R191$ & $1R93$\\
30 & $0R43$ & $1R44$ & 138 & $0R95$ & $1R156$\\
31 & $0R45$ & $1R46$ & 139 & $0R172$ & $1R163$\\
32 & $0R22$ & $1R47$ & 140 & $0R13$ & $1R156$\\
33 & $0R48$ & $1R33$ & 141 & $0R87$ & $1L163$\\
34 & $0R33$ & $1R34$ & 142 & $0R87$ & $1R155$\\
35 & $0L50$ & $1L3$ & 143 & $0L150$ & $1L160$\\
36 & $1R36$ & $1L2$ & 144 & $0R83$ & $1R87$\\
37 & $0R51$ & $1R37$ & 145 & $0R156$ & $1R13$\\
38 & $0R14$ & $1R53$ & 146 & $0R82$ & $1R141$\\
39 & $0L12$ & $0R37$ & 147 & $0R161$ & $1R165$\\
40 & $0L5$ & $1L54$ & 148 & $0L67$ & $1L139$\\
41 & $0R25$ & $0R26$ & 149 & $0R146$ & $1R199$\\
42 & $0R14$ & $0L17$ & 150 & $0R156$ & $1R113$\\
43 & $0R56$ & $1R57$ & 151 & $0R96$ & $1R96$\\
44 & $0R58$ & $1R59$ & 152 & $0R49$ & $0R33$\\
45 & $0R25$ & $0R55$ & 153 & $0R69$ & $1R200$\\
46 & $0R61$ & $1R62$ & 154 & $0L207$ & $1R95$\\
47 & $0R63$ & $0L17$ & 155 & $0R165$ & $1R176$\\
48 & $0R26$ & $1R48$ & 156 & $0R20$ & $1R167$\\
49 & $0L36$ & $0R64$ & 157 & $0R170$ & $1R162$\\
50 & $0L74$ & $1L3$ & 158 & $0R156$ & $1R95$\\
51 & $0R17$ & $1R51$ & 159 & $0R124$ & $1R161$\\
52 & $0R145$ & $1R141$ & 160 & $0R96$ & $1R159$\\
53 & $0R68$ & $1R69$ & 161 & $0L8$ & $1R140$\\
54 & $0L40$ & $1L54$ & 162 & $0L108$ & $1R170$\\
55 & $0R23$ & $1R55$ & 163 & $0R144$ & $1R169$\\
56 & $0L12$ & $0R51$ & 164 & $0R193$ & $1R29$\\
57 & $0R61$ & $0L17$ & 165 & $0R158$ & $1R141$\\
58 & $0R70$ & $1R71$ & 166 & $0R175$ & $1R175$\\
59 & $0R62$ & $1R72$ & 167 & $0R178$ & $1R121$\\
60 & $0R55$ & $1R60$ & 168 & $0R143$ & $1R162$\\
61 & $0R49$ & $0R26$ & 169 & $0R75$ & $1R61$\\
62 & $0L29$ & $1L73$ & 170 & $0R143$ & $1R87$\\
63 & $0L116$ & $1L31$ & 171 & $0R103$ & $1R95$\\
64 & $0R0$ & $1R64$ & 172 & $0R138$ & $1R83$\\
65 & $0R96$ & $1R148$ & 173 & $0R172$ & $1R75$\\
66 & $1R181$ & $1L54$ & 174 & $0R184$ & $0L204$\\
67 & $0R142$ & $1R143$ & 175 & $0R140$ & $1L147$\\
68 & $0R76$ & $1RH$ & 176 & $0R96$ & $1R52$\\
69 & $0L193$ & $1R32$ & 177 & $0R165$ & $1R52$\\
70 & $0R41$ & $1R63$ & 178 & $0R18$ & $1L130$\\
71 & $0L88$ & $1R77$ & 179 & $0R76$ & $1R115$\\
72 & $0R56$ & $1L78$ & 180 & $0R105$ & $1R123$\\
73 & $0R79$ & $1R80$ & 181 & $0L40$ & $1L109$\\
74 & $0R74$ & $1R81$ & 182 & $0R168$ & $1R156$\\
75 & $0L143$ & $1L139$ & 183 & $0R148$ & $1R114$\\
76 & $0R84$ & $1R85$ & 184 & $0R12$ & $1R97$\\
77 & $0R88$ & $1L131$ & 185 & $0R189$ & $1R98$\\
78 & $0L107$ & $1R57$ & 186 & $0R144$ & $1R195$\\
79 & $0R49$ & $0R34$ & 187 & $0R154$ & $1R188$\\
80 & $0R90$ & $0L42$ & 188 & $0R172$ & $1R182$\\
81 & $0R92$ & $1R194$ & 189 & $0R65$ & $1R8$\\
82 & $0R95$ & $1R13$ & 190 & $0R148$ & $1R126$\\
83 & $0R87$ & $1R96$ & 191 & $0R153$ & $1R111$\\
84 & $0R203$ & $1R22$ & 192 & $0R150$ & $1R122$\\
85 & $0R115$ & $1R99$ & 193 & $0R89$ & $1R87$\\
86 & $0R148$ & $1R112$ & 194 & $0R107$ & $1R214$\\
87 & $0R101$ & $1R102$ & 195 & $0R202$ & $1R108$\\
88 & $0R79$ & $1R78$ & 196 & $0R148$ & $1R166$\\
89 & $0R100$ & $1R171$ & 197 & $0R196$ & $1R157$\\
90 & $0R104$ & $1R105$ & 198 & $0R138$ & $1R136$\\
91 & $0R66$ & $0R1$ & 199 & $0R192$ & $1R19$\\
92 & $0R28$ & $1R106$ & 200 & $0R87$ & $1R205$\\
93 & $0R179$ & $0L204$ & 201 & $0R160$ & $1R190$\\
94 & $1L109$ & $1L94$ & 202 & $0R125$ & $1R149$\\
95 & $0R20$ & $1R110$ & 203 & $0R45$ & $1R99$\\
96 & $0R61$ & $1R61$ & 204 & $0L5$ & $0L204$\\
97 & $0R164$ & $1R154$ & 205 & $0R193$ & $1R138$\\
98 & $0R67$ & $1R186$ & 206 & $0R161$ & $1R207$\\
99 & $0R104$ & $0L204$ & 207 & $0R208$ & $1R179$\\
100 & $0R116$ & $1R117$ & 208 & $0R141$ & $1R209$\\
101 & $0R72$ & $1R59$ & 209 & $0R142$ & $1R210$\\
102 & $0R58$ & $1R118$ & 210 & $0R211$ & $1R86$\\
103 & $0R22$ & $1R119$ & 211 & $0R185$ & $1R162$\\
104 & $0R22$ & $0L11$ & 212 & $0R213$ & $1R38$\\
105 & $0R62$ & $1R42$ & 213 & $0R13$ & $1R38$\\
106 & $0R120$ & $1R9$ & 214 & $0R127$ & $1R212$\\
107 & $0R100$ & $1R121$ & & &\\
\bottomrule
\end{longtable}
\renewcommand{\arraystretch}{1}\normalsize

\end{document}
